% Shared preamble of the proof overview documents  <Folder>/<Name>_overview.tex.
% Each overview starts with
%     % Shared preamble of the proof overview documents  <Folder>/<Name>_overview.tex.
% Each overview starts with
%     % Shared preamble of the proof overview documents  <Folder>/<Name>_overview.tex.
% Each overview starts with
%     \input{../../scripts/overview_preamble.tex}
% (for ThinMatching/: \input{../scripts/overview_preamble.tex}) and is compiled from its own folder by
% scripts/build_overviews.py.
\documentclass[10pt,a4paper]{article}
\usepackage[T1]{fontenc}
\usepackage[utf8]{inputenc}
\usepackage[margin=2.2cm]{geometry}
\usepackage{amsmath,amssymb,amsthm}
\usepackage{array,booktabs,longtable}
\usepackage[dvipsnames]{xcolor}
\usepackage{enumitem}
\usepackage{tikz}
\usetikzlibrary{arrows.meta,positioning,shapes.multipart,fit,calc}
\usepackage[hidelinks]{hyperref}
\setlist{itemsep=2pt,topsep=4pt}
\renewcommand{\arraystretch}{1.2}

% Lean names, files and paper references
\newcommand{\lean}[1]{\textcolor{blue!50!black}{\texttt{\detokenize{#1}}}}
\newcommand{\file}[1]{\textcolor{OliveGreen!70!black}{\texttt{\detokenize{#1}}}}
\newcommand{\paperref}[1]{\textsf{#1}}
\newcommand{\hyp}[1]{\textcolor{BrickRed}{\textbf{#1}}}

% Boxes for the dependency diagrams
\tikzset{
  res/.style={draw, rounded corners, fill=blue!6, align=center, inner sep=4pt, font=\small},
  main/.style={res, fill=blue!18, very thick},
  cited/.style={res, fill=red!10, dashed},
  shared/.style={res, fill=gray!12},
  dep/.style={-{Stealth[length=2.2mm]}, thick},
}

\newcommand{\overviewtitle}[2]{%
  \begin{center}{\Large\bfseries #1}\\[4pt]{\small Proof overview of the Lean~4/Mathlib formalisation}\\[2pt]
  {\small #2}\end{center}}

% Standard paragraph describing the paper
\newcommand{\paperversion}{\noindent\textbf{Paper.} J.~Breitung, T.~Gervens, M.~Grohe.
  \emph{Graph Similarity Meets Optimal Transport}.\\ File \file{Graph_Similarity_Meets_Optimal_Transport.pdf} in the
  repository; statement, equation, page and line numbers refer to this PDF. \file{docs/NUMBERING.md} lists the Lean
  files and declarations of all numbered statements.\par\smallskip}

% (for ThinMatching/: \input{../scripts/overview_preamble.tex}) and is compiled from its own folder by
% scripts/build_overviews.py.
\documentclass[10pt,a4paper]{article}
\usepackage[T1]{fontenc}
\usepackage[utf8]{inputenc}
\usepackage[margin=2.2cm]{geometry}
\usepackage{amsmath,amssymb,amsthm}
\usepackage{array,booktabs,longtable}
\usepackage[dvipsnames]{xcolor}
\usepackage{enumitem}
\usepackage{tikz}
\usetikzlibrary{arrows.meta,positioning,shapes.multipart,fit,calc}
\usepackage[hidelinks]{hyperref}
\setlist{itemsep=2pt,topsep=4pt}
\renewcommand{\arraystretch}{1.2}

% Lean names, files and paper references
\newcommand{\lean}[1]{\textcolor{blue!50!black}{\texttt{\detokenize{#1}}}}
\newcommand{\file}[1]{\textcolor{OliveGreen!70!black}{\texttt{\detokenize{#1}}}}
\newcommand{\paperref}[1]{\textsf{#1}}
\newcommand{\hyp}[1]{\textcolor{BrickRed}{\textbf{#1}}}

% Boxes for the dependency diagrams
\tikzset{
  res/.style={draw, rounded corners, fill=blue!6, align=center, inner sep=4pt, font=\small},
  main/.style={res, fill=blue!18, very thick},
  cited/.style={res, fill=red!10, dashed},
  shared/.style={res, fill=gray!12},
  dep/.style={-{Stealth[length=2.2mm]}, thick},
}

\newcommand{\overviewtitle}[2]{%
  \begin{center}{\Large\bfseries #1}\\[4pt]{\small Proof overview of the Lean~4/Mathlib formalisation}\\[2pt]
  {\small #2}\end{center}}

% Standard paragraph describing the paper
\newcommand{\paperversion}{\noindent\textbf{Paper.} J.~Breitung, T.~Gervens, M.~Grohe.
  \emph{Graph Similarity Meets Optimal Transport}.\\ File \file{Graph_Similarity_Meets_Optimal_Transport.pdf} in the
  repository; statement, equation, page and line numbers refer to this PDF. \file{docs/NUMBERING.md} lists the Lean
  files and declarations of all numbered statements.\par\smallskip}

% (for ThinMatching/: \input{../scripts/overview_preamble.tex}) and is compiled from its own folder by
% scripts/build_overviews.py.
\documentclass[10pt,a4paper]{article}
\usepackage[T1]{fontenc}
\usepackage[utf8]{inputenc}
\usepackage[margin=2.2cm]{geometry}
\usepackage{amsmath,amssymb,amsthm}
\usepackage{array,booktabs,longtable}
\usepackage[dvipsnames]{xcolor}
\usepackage{enumitem}
\usepackage{tikz}
\usetikzlibrary{arrows.meta,positioning,shapes.multipart,fit,calc}
\usepackage[hidelinks]{hyperref}
\setlist{itemsep=2pt,topsep=4pt}
\renewcommand{\arraystretch}{1.2}

% Lean names, files and paper references
\newcommand{\lean}[1]{\textcolor{blue!50!black}{\texttt{\detokenize{#1}}}}
\newcommand{\file}[1]{\textcolor{OliveGreen!70!black}{\texttt{\detokenize{#1}}}}
\newcommand{\paperref}[1]{\textsf{#1}}
\newcommand{\hyp}[1]{\textcolor{BrickRed}{\textbf{#1}}}

% Boxes for the dependency diagrams
\tikzset{
  res/.style={draw, rounded corners, fill=blue!6, align=center, inner sep=4pt, font=\small},
  main/.style={res, fill=blue!18, very thick},
  cited/.style={res, fill=red!10, dashed},
  shared/.style={res, fill=gray!12},
  dep/.style={-{Stealth[length=2.2mm]}, thick},
}

\newcommand{\overviewtitle}[2]{%
  \begin{center}{\Large\bfseries #1}\\[4pt]{\small Proof overview of the Lean~4/Mathlib formalisation}\\[2pt]
  {\small #2}\end{center}}

% Standard paragraph describing the paper
\newcommand{\paperversion}{\noindent\textbf{Paper.} J.~Breitung, T.~Gervens, M.~Grohe.
  \emph{Graph Similarity Meets Optimal Transport}.\\ File \file{Graph_Similarity_Meets_Optimal_Transport.pdf} in the
  repository; statement, equation, page and line numbers refer to this PDF. \file{docs/NUMBERING.md} lists the Lean
  files and declarations of all numbered statements.\par\smallskip}

\begin{document}
\overviewtitle{Proposition 13: a large integrality gap for the spectral distance}{Folder \file{GsimOT/Proposition14_SpectralLowerBound}}
\paperversion

\section{Statement}

\paragraph{Paper.} Proposition~13 (Section~3.4, p.~11, l.~389--390): \emph{for every $r\ge1$ there are graphs
$G_r,H_r$ on $n=\Theta(r^2)$ vertices such that $d_2(G_r,H_r)=\sqrt r$ and
$\frac1{2\sqrt2}\le d^{\mathrm{OT}}_2(G_r,H_r)\le1$.} The text after it (l.~391--393) concludes that there is no
uniform bound $\delta_2\le C_\alpha(\delta^{\mathrm{OT}}_2)^\alpha$ for $\alpha>\frac34$. Theorem~3~(4) (p.~4,
l.~127): $\delta_2(G_r,H_r)\ge\sqrt r\cdot\delta^{\mathrm{OT}}_2(G_r,H_r)$. Proof: App.~F.1 (pp.~45--46), with the
construction, Figure~7 and Lemma~49 ($\frac r{\sqrt n}\le d^{\mathrm{OT}}_2(G_r,H_r)\le1$).

\paragraph{Lean.} The graphs are \lean{GsimOT.Prop14.Gr r}, \lean{GsimOT.Prop14.Hr r} on \lean{Vtx r}. For every
\hyp{$r\ge1$}:
\begin{itemize}
\item \lean{proposition13}: $|V|=4r^2+3r+1$, $d_2(G_r,H_r)=\sqrt r$ and $\frac1{2\sqrt2}\le d^{\mathrm{OT}}_2(G_r,H_r)\le1$
  (exactly the statement);
\item \lean{proposition14}: Proposition~13 together with the bound $\frac r{\sqrt n}\le d^{\mathrm{OT}}_2(G_r,H_r)$
  of Lemma~49, i.e.\ the same with $\frac1{2\sqrt2}\le\frac r{\sqrt n}\le d^{\mathrm{OT}}_2(G_r,H_r)$;
  \lean{lemma53}: Lemma~49;
\item \lean{theorem4_item4}: Theorem~3~(4);
\item \lean{no_power_bound}: for all $\alpha>\frac34$ and $C\in\mathbb R$ there is $r\ge1$ with
  $C\,\delta^{\mathrm{OT}}_2(G_r,H_r)^\alpha<\delta_2(G_r,H_r)$.
\end{itemize}

\section{Where the proof lives}
\begin{longtable}{>{\raggedright\arraybackslash}p{5.2cm} >{\raggedright\arraybackslash}p{10.6cm}}
\toprule \textbf{File} & \textbf{Content} \\ \midrule\endhead \bottomrule\endfoot
\file{Construction.lean} & the vertex set $X\cup M$ (\lean{Core}, \lean{Marker}, \lean{Vtx}), the graphs
  (\lean{graph}, \lean{Gr}, \lean{Hr}), $|V|=4r^2+3r+1$ (\lean{card_Vtx}), degrees (\lean{le_deg_Hr}). \\
\file{SpectralDistance.lean} & $d_2(G_r,H_r)=\sqrt r$ (\lean{distPQ_eq}) and elementary facts on $\ell^2$ and
  spectral norms. \\
\file{Lemma53.lean} & Lemma~49: upper bound via $Q=\frac1{2r}\sum_iP_{\sigma_i}$ (\lean{dOT_le_one}), lower bound via
  the test vector $\mathbf 1_R$ (\lean{le_dOT_of_deg}, \lean{le_dOT}). \\
\file{Proposition14.lean} & Proposition~13 (\lean{proposition13}, \lean{proposition14}), Theorem~3~(4), and the
  remark (\lean{no_power_bound}). \\
\file{Common/OTDistances.lean}, \file{Common/OperatorNorm.lean} & $B_{2\to2}$, $d^{\mathrm{OT}}_2$,
  $\delta^{\mathrm{OT}}_2=n^{-1}d^{\mathrm{OT}}_2$ (\lean{deltaOTPQ_eq_dOT}), Lemma~19. \\
\end{longtable}

\section{Proof layout}
\begin{center}
\begin{tikzpicture}
  \node[main] (p13) at (-3.5,0) {Proposition 13\\\lean{proposition13}};
  \node[main] (t3) at (1.5,0) {Theorem 3 (4)\\\lean{theorem4_item4}};
  \node[main] (rem) at (6,0) {remark, $\alpha>3/4$\\\lean{no_power_bound}};
  \node[main] (p14) at (-3.5,-2) {Prop.\ 13 + Lemma 49\\\lean{proposition14}};
  \node[res] (norm) at (3.5,-2) {normalised: $\delta_2=\sqrt r/n$, $\delta^{\mathrm{OT}}_2\le1/n$\\\lean{deltaPQ_eq}, \lean{deltaOTPQ_bounds}};
  \node[res] (d2) at (-6,-4.3) {$d_2(G_r,H_r)=\sqrt r$\\\lean{distPQ_eq}};
  \node[res] (l49) at (-1,-4.3) {Lemma 49\\\lean{lemma53}};
  \node[shared] (rel) at (5,-4.3) {$\delta^{\mathrm{OT}}_2=n^{-1}d^{\mathrm{OT}}_2$\\\lean{deltaOTPQ_eq_dOT}};
  \node[res] (up) at (-3.8,-6.6) {upper bound $C_2(Q)\le1$\\\lean{opNorm_Bq_le}, \lean{dOT_le_one}};
  \node[res] (low) at (2.2,-6.6) {lower bound $r/\sqrt n$\\\lean{le_opNorm_Bmat}, \lean{le_dOT}};
  \node[res] (con) at (-6.2,-8.9) {construction\\\lean{Gr}, \lean{Hr}, \lean{card_Vtx}, \lean{le_deg_Hr}};
  \node[shared] (l19) at (0,-8.9) {Lemma 19\\\lean{exists_abs_dotProduct_mulVec_eq_opNorm}};
  \draw[dep] (p13) -- (p14); \draw[dep] (t3) -- (norm); \draw[dep] (rem) -- (norm);
  \draw[dep] (p14) -- (d2); \draw[dep] (p14) -- (l49);
  \draw[dep] (norm) -- (l49); \draw[dep] (norm) -- (d2); \draw[dep] (norm) -- (rel);
  \draw[dep] (l49) -- (up); \draw[dep] (l49) -- (low);
  \draw[dep] (up) -- (l19); \draw[dep] (up) -- (con); \draw[dep] (d2) -- (con);
\end{tikzpicture}
\end{center}
Arrows mean ``uses''. Blue: this folder; grey: \file{GsimOT/Common}. Nothing is cited.

\section{The proof step by step}

\paragraph{Construction (p.~45, l.~1399--1411; Figure~7, p.~45).} Core $X=\{u\}\cup\{v_{i,k}\mid i\in[2r],0\le k\le r\}$
is \lean{Core r} $=$ \texttt{Option (Fin (2r) $\times$ Fin (r+1))} ($u=$ \texttt{none}); markers
$M=C\cup T_1\cup\dots\cup T_{2r}$ are \lean{Marker r} (clocks \texttt{inl j} $=c_{j+1}$, tags \texttt{inr (i,t)}).
$N(v_{i,k})=T_i\cup\{c_1,\dots,c_k\}$ in both graphs, $u$ is isolated in $G_r$ and adjacent to all clocks in $H_r$:
one definition \lean{graph} with a Boolean flag. $|V|=4r^2+3r+1$ is \lean{card_Vtx}; $H_r$ has minimum degree
$\ge r$ (\lean{le_deg_Hr}).

\paragraph{$d_2(G_r,H_r)=\sqrt r$ (p.~45, l.~1412--1416).} Lower bound: the spectral norm is at least the Euclidean
norm of a row (\lean{sqrt_row_le_opNorm22}), and the row of $\pi(u)$ has $\ge r$ entries $\pm1$
(\lean{sqrt_le_opNorm_alignedDiff}). Upper bound: for the identity, $A_G-A_H$ is minus a star with $r$ leaves, of
norm $\sqrt r$ (\lean{opNorm_diff_refl_le}). Together: \lean{distPQ_eq}.

\paragraph{Lemma 49, upper bound (pp.~45--46, l.~1418--1436).} $\sigma_i$ shifts branch $i$ one step along
$u\to v_{i,0}\to\dots\to v_{i,r}\to u$ (\lean{shift}, \lean{tau}); $Q=\frac1{2r}\sum_iP_{\sigma_i}$ fixes the markers
(\lean{Qs}, \lean{Qs_isDS}). The vectors $d_{xy}=\mathbf 1_{N_{G}(x)}-\mathbf 1_{N_H(y)}$ (\lean{dv}) are
$-\mathbf 1_{T_i}$, $-\mathbf 1_{\{c_{k+1}\}}$, $\mathbf 1_{T_i}$ on the support. $B_{2\to2}(G_r,H_r,Q)$ has the block
form $\left(\begin{smallmatrix}0&K\\K^\top&0\end{smallmatrix}\right)$ (\lean{Bq_mulVec_cc}, \lean{Bq_mulVec_mm},
\lean{dotProduct_Bq}). The paper computes $\|K^\top K\|_2=\|\frac1r\sum_i\mathbf 1_{T_i}\mathbf 1_{T_i}^\top+\sum_j
\mathbf 1_{c_j}\mathbf 1_{c_j}^\top\|_2=1$. The Lean proves the corresponding quadratic-form bound
$\zeta^\top K^\top K\zeta=\sum_{x,y}Q(x,y)(d_{xy}\cdot\zeta)^2\le\|\zeta\|^2$ (\lean{sum_branch_le},
\lean{sum_qc_dvDot_le}: Cauchy--Schwarz on each tag block, orthogonality of the clocks), and concludes
$|y^\top Bx|\le1$ for unit $x,y$ via Lemma~19 (\lean{opNorm_Bq_le}), hence $d^{\mathrm{OT}}_2\le1$
(\lean{dOT_le_one}).

\paragraph{Lemma 49, lower bound (p.~46, l.~1437--1446).} For every $R\in\mathrm{DS}$, the unit vector
$\mathbf 1_R(x,y)=\sqrt{R(x,y)/n}$ satisfies $(B\mathbf 1_R)(x,y)=\mathbf 1_R(x,y)(\deg_G x-\deg_H y)$, and the rows
of $u$ give $\|B\mathbf 1_R\|\ge r/\sqrt n$. This is proved for arbitrary graphs and any vertex $u$ with
$|\deg_G u-\deg_H y|\ge c$ for all $y$ (\lean{le_opNorm_Bmat}, \lean{le_dOT_of_deg}), then specialised
(\lean{le_dOT}).

\paragraph{Proof of Proposition 13 (p.~46, l.~1447--1448).} Lemma~49, $d_2=\sqrt r$ and $\frac r{\sqrt n}\ge
\frac1{2\sqrt2}$ since $8r^2\ge n$: \lean{proposition14}; \lean{proposition13} drops the middle term.

\paragraph{Theorem 3 (4) and the remark (p.~4, l.~127; p.~11, l.~391--393).} $\delta_2(G_r,H_r)=\sqrt r/n$
(\lean{deltaPQ_eq}) and $0\le\delta^{\mathrm{OT}}_2(G_r,H_r)\le1/n$ (\lean{deltaOTPQ_bounds}, via
\lean{deltaOTPQ_eq_dOT}) give \lean{theorem4_item4}. For \lean{no_power_bound}: for $\alpha>\frac34$ choose
$a=\min\{\alpha,1\}$ and $r$ large; $(\delta^{\mathrm{OT}}_2)^\alpha\le n^{-a}$ and $n\le8r^2$ give
$C(\delta^{\mathrm{OT}}_2)^\alpha\le C8^{1-a}r^{2(1-a)}/n<\sqrt r/n$. Only the upper bound
$\delta^{\mathrm{OT}}_2=O(n^{-1})$ is needed, not $\Theta(n^{-1})$.

\section{Assumptions and trusted base}
No cited results; the hypothesis is $r\ge1$. Axioms: \texttt{propext}, \texttt{Classical.choice},
\texttt{Quot.sound}.

\section{Notes on the formalisation}
\begin{itemize}
\item The upper bound of Lemma~49 is proved through the quadratic form of $K^\top K$, and only $C_2(Q)\le1$ is
  proved, which is what $d^{\mathrm{OT}}_2(G_r,H_r)\le1$ requires (the paper computes $C_2(Q)=1$).
\item The lower bound of Lemma~49 is proved for arbitrary graphs.
\item \lean{proposition14} records the intermediate bound $\frac r{\sqrt n}\le d^{\mathrm{OT}}_2(G_r,H_r)$ used in the
  proof of Proposition~13; \lean{proposition13} is derived from it.
\end{itemize}

\section{How to check this result}
\begin{verbatim}
lake build GsimOT.Proposition14_SpectralLowerBound.Proposition14
lake env lean scripts/CheckAxioms.lean
\end{verbatim}
The table in \file{Proposition14_correspondence.pdf} lists every step with file and line.
\end{document}
